% ============================================================================
\chapter{S\'ecurit\'e web : ce qu'on doit savoir avant de mettre en ligne}
\label{ch:js-secu}
% ============================================================================

\begin{objectif}
\begin{itemize}
  \item Int\'egrer une r\`egle unique~: \textbf{toute donn\'ee qui vient de l'ext\'erieur est hostile}.
  \item Comprendre et bloquer les six attaques qui font 90 \% des incidents~: XSS, CSRF,
        injection SQL, IDOR, upload malveillant, d\'ependance compromise.
  \item Configurer en-t\^etes, cookies, CSP et politique de session correctement.
  \item Hacher un mot de passe comme il faut, et ne jamais \og ~inventer de la cryptographie~ \fg{}.
  \item Savoir o\`u chercher quand on veut aller plus loin (OWASP, outils, pratiques).
\end{itemize}
\end{objectif}

\section{Les trois principes qui rendent tout le reste \'evident}

\begin{enumerate}
  \item \textbf{Sortie encod\'ee, pas entr\'ee filtr\'ee.} Ne cherche pas \`a nettoyer une saisie pour
        qu'elle soit \og ~s\^ure~ \fg{}~: encode-la \emph{au moment de l'afficher}, selon le contexte
        (HTML, attribut, URL, JS, SQL). Un \verb|<script>alert(1)</script>| dans un champ \og ~nom~ \fg{}
        est une donn\'ee valide~; c'est son rendu qui doit devenir inoffensif.
  \item \textbf{D\'enominateur de confiance.} Le navigateur ne sait pas qui tu es~: il ne peut pas
        authentifier un appel. La confiance vient d'un secret partag\'e (cookie \verb|HttpOnly|,
        jeton sign\'e) \emph{et} d'une v\'erification d'autorisation \`a chaque requ\^ete.
  \item \textbf{Moindre privil\`ege.} Un utilisateur ne voit que ses donn\'ees. Un compte de base de
        donn\'ees n'a pas \verb|DROP|. Un conteneur ne tourne pas en root. Un jeton d'API a une
        port\'ee et une dur\'ee.
\end{enumerate}

\section{XSS~: le script malveillant dans ta page}

Un XSS, c'est du code que l'attaquant fait ex\'ecuter \emph{avec tes droits} dans le navigateur de
la victime~: vol de session, usurpation d'actions, enregistrement des frappes.

\begin{codejs}
// VULNERABLE : du texte utilisateur rendu comme du HTML
el.innerHTML = commentaire.texte;                      // <img src=x onerror="steal(document.cookie)">
document.querySelector("#u").innerHTML = `<a href="${utilisateur.site}">site</a>`;  // href="javascript:..."

// SUR : le texte reste du texte
el.textContent = commentaire.texte;
a.href = new URL(utilisateur.site, location.href).protocol.startsWith("http") ? utilisateur.site : "#";

// SI tu dois absolument rendre du HTML (markdown, richtext) : passer par un assainisseur
import DOMPurify from "dompurify";
el.innerHTML = DOMPurify.sanitize(markdownVersHtml(commentaire.texte), {
  ALLOWED_TAGS: ["p", "strong", "em", "a", "ul", "ol", "li", "code", "pre", "blockquote"],
  ALLOWED_ATTR: ["href", "title", "class"],
  ALLOW_DATA_ATTR: false,
  ADD_FORBID_TAGS: ["style"],
});
// Cote serveur, pour du contenu stocke une fois pour toutes : sanitize-html, DOMPurify,
// ou DOMPurify en JSDOM. Ne stocke JAMAIS du HTML non filtre.
\end{codejs}

\etape{Content Security Policy~: le filet de s\'ecurit\'e}
La CSP ne remplace pas l'encodage~; elle rend l'exploitation tr\`es difficile en interdisant
le JS inline et les sources inconnues.

\begin{codeconf}
# En-tete envoye par ton serveur (Express : helmet l'fait pour toi, personnalise-le)
Content-Security-Policy: default-src 'self'; base-uri 'self'; object-src 'none';
  frame-ancestors 'none'; form-action 'self'; upgrade-insecure-requests;
  script-src 'self' https://cdn.example.com 'nonce-{NONCE}';
  style-src 'self' 'unsafe-inline'; img-src 'self' data: https:; font-src 'self';
  connect-src 'self' https://api.nour.dev; report-uri /csp-report
Report-To: {"group":"csp","max_age":10886400,"endpoints":[{"url":"https://csp.nour.dev/reports"}]}
Cross-Origin-Opener-Policy: same-origin
Cross-Origin-Embedder-Policy: require-corp
Cross-Origin-Resource-Policy: same-origin
\end{codeconf}

Avec une nonce~: \verb|<script nonce="{valeur-aleatoire-par-requete}">|~; le JS inline sans nonce
est bloqu\'e. \'Evite \verb|unsafe-inline| et \verb|unsafe-eval|~: si tu les mets, la CSP est
d\'ecorative. Pour les frameworks qui exigent du style inline, utilise
\verb|style-src 'self' 'unsafe-inline'| (moins grave) ou un hachage.

\begin{itemize}
  \item G\'en\'erateur et validateur~: \link{CSP~Evaluator}{https://csp-evaluator.withgoogle.com/}.
  \item \texttt{Trusted Types} (Chrome/Edge/Firefox r\'ecents) rend \texttt{innerHTML} refus\'e sans
        une politique explicite~: \url{https://developer.mozilla.org/en-US/docs/Web/API/Trusted\\\\\_Types\\\\\_API}.
  \item \texttt{SRI} pour les scripts tiers~: \verb|<script src="..." integrity="sha384-..." crossorigin>|.
\end{itemize}

\section{CSRF~: le site d'un autre fait des requ\^etes en ton nom}

Un \verb|POST /compte/modifier-email| accept\'e sur la seule base du cookie de session est
exploitable depuis n'importe quelle page web. Trois d\'efenses, cumulables~:

\begin{enumerate}
  \item Cookie \verb|SameSite=Lax| (d\'efaut des navigateurs modernes)~: bloque la plupart des cas.
  \item V\'erifier l'\verb|Origin|/\verb|Referer| de la requ\^ete mutante~: peu co\^uteux, tr\`es efficace.
  \item Jeton CSRF~: une valeur secr\`ete al\'eatoire par session, renvoy\'ee dans le formulaire ou
        un en-t\^ete \verb|X-CSRF-Token|, recontr\^ol\'ee c\^ot\'e serveur pour toute m\'ethode
        non s\^ure (POST/PUT/PATCH/DELETE).
\end{enumerate}

\begin{codejs}
// Generation (cote serveur) puis controle. Le jeton est lie a la session, pas au navigateur.
import { randomBytes, timingSafeEqual } from "node:crypto";

function nouveauJeton(req) {
  req.session.csrf ??= randomBytes(32).toString("base64url");
  return req.session.csrf;
}
function verifierJeton(req) {
  const recu = req.get("x-csrf-token") ?? req.body?._csrf;
  const attendu = req.session?.csrf;
  if (!attendu || !recu) return false;
  const a = Buffer.from(recu), b = Buffer.from(attendu);
  return a.length === b.length && timingSafeEqual(a, b);       // comparaison temps constant
}
app.use((req, res, next) => {
  if (["GET", "HEAD", "OPTIONS"].includes(req.method)) return next();
  if (!verifierJeton(req)) return res.status(403).json({ erreur: "jeton CSRF manquant ou invalide" });
  next();
});
\end{codejs}

\section{Injection---le classique qui n'a pas vieilli}

\begin{codesql}
-- Ce que produit une concatenation naive :
SELECT * FROM users WHERE nom = '' OR '1'='1' -- '   -> TOUS les utilisateurs
-- En PHP comme en Node comme en Java : REQUETE PREPAREE, TOUJOURS.
\end{codesql}

\begin{codejs}
// Node + Postgres : les placeholders font le travail
const { rows } = await pool.query("SELECT * FROM articles WHERE slug = $1 AND publie = true", [slug]);

// Morceau de requete dynamique (tri, colonnes) : les placeholders ne suffisent pas.
// On ne fait jamais confiance a la chaine : on la CONFRONTE a une liste blanche.
const COLONNES = { recent: "cree_le", prix: "prix", nom: "nom" };
const ordre = COLONNES[req.query.tri] ?? "cree_le";
const sens = String(req.query.ordre).toLowerCase() === "asc" ? "ASC" : "DESC";
await pool.query(`SELECT * FROM articles ORDER BY ${ordre} ${sens} LIMIT $1 OFFSET $2`, [taille, decale]);
\end{codejs}

Le m\^eme principe prot\`ege des injections dans un shell~: ne jamais construire une commande avec
de la concat\'enation~; passe un \emph{tableau d'arguments} (\verb|execFile("git", [args])|,
\verb|execFile| et non \verb|exec|). Et pour les noms de fichiers~: \verb|path.resolve| puis
v\'erifier que le r\'esultat commence par le dossier autoris\'e.

\section{IDOR, autorisation, et le pi\`ege du \og ~c'est mon id~ \fg{}}

\begin{codejs}
// VULNERABLE : l'identifiant fourni par le client designe la ressource, sans controle
app.get("/api/factures/:id", async (req, res) => res.json(await facture(req.params.id)));
// -> n'importe qui remplace :id par celui du voisin.

// CORRECT : la requete est FILTREE par le proprietaire authentifie
app.get("/api/factures/:id", async (req, res) => {
  const f = await facturePour(req.user.id, req.params.id);   // WHERE id = $1 AND utilisateur_id = $2
  if (!f) return res.status(404).json({ erreur: "introuvable" });   // pas 403 : ne pas fuiter l'existence
  res.json(f);
});

// Autorisation centralisee, pas diss\'eminee dans les controleurs
const politiques = {
  "article:lire": (u, a) => a.publie || u.id === a.auteurId || u.roles.includes("editeur"),
  "article:ecrire": (u, a) => u.id === a.auteurId || u.roles.includes("editeur"),
  "compte:modifier": (u) => true,
};
export function autoriser(u, action, ressource) {
  const p = politiques[action];
  if (!p) throw new Error(`action inconnue : ${action}`);
  return Boolean(p(u, ressource));
}
\end{codejs}

\section{Mots de passe, sessions, jetons}

\begin{codejs}
// Hachage : argon2id (le meilleur) ou bcrypt (le plus courant). Jamais sha256(mdpmot + sel).
import argon2 from "argon2";
const hash = await argon2.hash(mdp, { type: argon2.argon2id, memoryCost: 19_456, timeCost: 2, parallelism: 1 });
await argon2.verify(hash, mdpSoumis);        // temps constant, sel int\'egr\'e

// Politique raisonnable (NIST SP 800-63B) : longueur >= 12, pas de composition absurde,
// comparaison aux mots de passe fuit\'es, pas de date d'expiration arbitraire.
import { createHash } from "node:crypto";
const hain = createHash("sha1").update(mdp).digest("hex");
const [prefixe, ...reste] = [hain.slice(0, 5), hain.slice(5)];
const rep = await fetch(`https://api.pwnedpasswords.com/range/${prefixe}`);
if ((await rep.text()).split("\n").some((l) => l.startsWith(reste.join("")))) {
  throw new Error("Ce mot de passe figure dans des fuites connues : choisis-en un autre");
}

// Reinitialisation : jeton a usage unique, aleatoire, court (15 min), stocke HACHE en base,
// envoye par email, invalide la session precedente, et reponse identique que l'email existe ou non.
\end{codejs}

\begin{tabz}{lX}
  \headrow \thd{Sujet} \& \thd{R\'eglage s\^ur} \\
  \midrule
  Cookie de session \& \verb|HttpOnly; Secure; SameSite=Lax; Path=/|, expiration 30 min -- 12 h. \\
  Rotation \& r\'eg\'en\'erer l'identifiant de session \`a la connexion et \`a l'\'el\'evation de privil\`ege. \\
  D\'econnexion \& invalider c\^ot\'e serveur (pas seulement effacer le cookie). \\
  \og ~Rester connect\'e~ \fg{} \& cookie longue dur\'ee avec un jeton \emph{stock\'e c\^ot\'e serveur}, r\'eversible. \\
  JWT \& courte dur\'ee (5--15 min), \verb|alg| \'epingl\'e (refuser \verb|none|), pas de donn\'ees sensibles, r\'evocation par liste ou refresh en cookie. \\
  Cl\'e API \& port\'ee minimale, rotation, jamais dans le code ni dans un bundle. \\
  CSRF \& jeton ou \verb|SameSite| + v\'erification d'\verb|Origin|. \\
  2FA \& TOTP \link{otplib}{https://github.com/yeqown/otplib} ou WebAuthn (bien meilleur) --- \url{https://webauthn.guide/} \\
\end{tabz}

\section{En-t\^etes de base, \`a poser partout}

\begin{codeconf}
Strict-Transport-Security: max-age=63072000; includeSubDomains; preload
X-Content-Type-Options: nosniff
X-Frame-Options: DENY                     (redondant avec frame-ancestors 'none', encore utile)
Referrer-Policy: strict-origin-when-cross-origin
Permissions-Policy: camera=(), microphone=(), geolocation=(), interest-cohort=()
Content-Security-Policy: (voir plus haut)
Server: (retirer ou rendre vague)
X-Powered-By: (retirer)
\end{codeconf}

\section{Upload, liens, et les petits d\'etails qui tuent}

\begin{itemize}
  \item \textbf{Upload}~: limite la taille, liste blanche d'\emph{extensions ET} de types MIME
        (v\'erifi\'es par analyse du contenu, pas par le nom), stocke hors racine web avec un nom
        g\'en\'er\'e, sert avec \verb|content-disposition: attachment| pour le HTML/svg, et jamais
        ex\'ecutable. Un \verb|.svg| peut contenir un \verb|<script>|~; un \verb|.php.jpg| peut
        tromper une v\'erification na\"ive.
  \item \textbf{Redirections ouvertes}~: \verb|?next=//evil.com|~: n'accepte que des chemins
        relatifs ou une liste blanche d'h\^otes.
  \item \textbf{Regex et ReDoS}~: limite la longueur avant de matcher.
  \item \textbf{D\'eni de service par requ\^ete co\^uteuse}~: paginer obligatoirement, limiter la
        taille des listes, mettre un cache.
  \item \textbf{Logs}~: ne jamais journaliser mot de passe, jeton, cookie, carte, document
        d'identit\'e. Masque et tronque.
  \item \textbf{Erreurs}~: en production, aucune stack trace, aucun SQL, aucun chemin de fichier
        au visiteur. (Les logs c\^ot\'e serveur, eux, gardent tout.)
  \item \textbf{D\'ependances}~: \verb|npm audit|, \link{Socket}{https://socket.dev/},
        lockfile verrouill\'e, \link{Renovate}{https://github.com/renovatebot/renovate} ou
        Dependabot pour les mises \`a jour automatiques.
\end{itemize}

\begin{terminal}
# Auto-\'evaluation d'un site que TU poss\`edes (ou un localhost que tu montes)
npx @lhce/lighthouse-cli --only-categories=security http://localhost:4173
npm i -D zap-baseline && npx zap-baseline.py -t http://localhost:3000 -r rapport.html
# ZAP (benchmarks OWASP) : https://www.zaproxy.org/   ---   Nuclei : https://github.com/projectdiscovery/nuclei
# Test manuel des en-tetes : https://securityheaders.com/  (public)
\end{terminal}

\begin{exercice}
\begin{enumerate}
  \item Fais tomber toi-m\^eme ton application de la semaine pr\'ec\'edente~: XSS stock\'e dans un
        commentaire, XSS par URL dans un lien, redirection ouverte, IDOR sur les factures,
        CSRF sur le changement d'email. \emph{Note l'exploit exact}, puis corrige-le, puis
        rejoue-le pour v\'erifier.
  \item Ajoute une CSP fonctionnelle en mode \verb|Content-Security-Policy-Report-Only|,
        corrige toutes les violations remont\'ees, puis active-la vraiment.
  \item \'Ecris un test automatis\'e pour chaque faille corrig\'ee (un test = une vuln\'erabilit\'e
        referm\'ee pour toujours).
  \item Audite ton \verb|package.json|~: nombre de d\'ependances transitives, licences, \^age des
        versions, paquets inutilis\'es~; remplace ce qui doit l'\^etre.
  \item Rends ta page d'authentification conforme~: cookie \verb|HttpOnly|, rotation de session,
        limitation de d\'ebit (5 essais/minute/IP), message identique que l'email existe ou non.
  \item Lis le OWASP Top~10 en une heure et, pour chaque point, \'ecris une ligne sur la fa\c{c}on
        dont ton projet actuel s'en prot\`ege ou ne s'en prot\`ege pas.
\end{enumerate}
\end{exercice}

\begin{ressources}
\begin{itemize}
  \item OWASP Top~10 (la r\'ef\'erence) --- \url{https://owasp.org/Top10/fr/}
  \item OWASP Cheat Sheet Series (fiches courtes et pr\'ecises) --- \url{https://cheatsheetseries.owasp.org/}
  \item MDN S\'ecurit\'e web --- \url{https://developer.mozilla.org/fr/docs/Web/Security}
  \item PortSwigger Web Security Academy (gratuit, exercices interactifs, le meilleur cours pratique) --- \url{https://portswigger.net/web-security}
  \item OWASP ZAP --- \url{https://www.zaproxy.org/} --- securityheaders.com
  \item \url{https://haveibeenpwned.com/} --- v\'erifier une fuite de son propre email.
\end{itemize}
\end{ressources}

\index{securite}\index{XSS}\index{CSRF}\index{CSP}
